\subsection*{\texorpdfstring{$\lambda$-Calculus}{-Calculus}}

\kwub{Lambda expressions} are composed of:
\begin{itemize}
    \vItem
    \kwb{variables} \iMbox{v_1,v_2, \dots};
    \vItem
    \kwb{abstraction symbols} \iMbox{\lambda} (lambda) and \iMbox{\left. . \right.} (dot);
    \vItem
    \kwb{parentheses} \iMbox{()}
\end{itemize}

Set of \kwub{lambda expressions} \iMbox{\Lambda} => \stu{defined inductively}:
\begin{enumerate}
    \vItem \kwb{Variable} => \iMbox{x \ \text{is a variable} \implies x\in \Lambda};
    \vItem \kwb{Abstraction} =>
    \iMbox{ x \ \text{is a variable and} \ M\in \Lambda \implies ( \lambda x . M) \in \Lambda};
    \vItem \kwb{Application} =>
    \iMbox{ M,N \in \Lambda \implies ( M \ N ) \in \Lambda};

    \vItem Or in \stu{Backus-Naur form}:
    \iMbox{\begin{aligned}
             & <expression>  & ::= \ \  & <abstraction>                 \\
             &               & | \ \    & <application>                 \\
             &               & | \ \    & <variable>                    \\
             & <abstraction> & ::= \ \  & λ <variable> . <expression>   \\
             & <application> & ::= \ \  & ( <expression> <expression> ) \\
             & <variable>    & ::= \ \  & v1 | v2 | ...
        \end{aligned}}
\end{enumerate}

\kwub{Conventional Notation}
\begin{itemize}
    \vItem
    \stu{Outermost parentheses} dropped: \iMbox{M\ N} instead of \iMbox{(M\ N)}
    \vItem \kwb{left associative applications}: \iMbox{M\ N \ P= ((M \ N) \ P)}
    \vItem \stu{Space in applications omitted} for \stu{single-letter} variables: \iMbox{xyz = x \ y \ z}
    \vItem \stu{Lambda abstraction} has \stu{lower precedence} than
    \stu{application}: \iMbox{\lambda x.M \ N = \lambda x. (M \ N)}
    \vItem \stu{Sequence of abstractions} is \kwb{contracted}: \iMbox{\lambda x. \lambda y. \lambda z. M = \lambda xyz. M}
\end{itemize}

\hSep

\kwub{Free and bound variables}
\begin{enumerate}
    \vItem The abstraction operator \iMbox{\lambda} \stb{binds its variable}
    wherever it occurs \stu{in the body} of the abstraction
    \vItem Variables that fall \stu{within the scope} of an abstraction are said to
    be \kwb{bound}
    \vItem All other variables are called \kwb{free}
    \vItem \textbf{NOTE:} A variable is \stu{bound by its nearest abstraction}
    \begin{enumerate}
        \vItem
        e.g. in \iMbox{\lambda y. x x y}, the \iMbox{y} is
        \stu{bound} and the \iMbox{x} is \stu{free}
        \vItem
        e.g. in \iMbox{\lambda x. y (\lambda x. xz)} the \iMbox{x} is
        \stu{bound} to the second lambda
    \end{enumerate}
\end{enumerate}

\hSep

\stub{Set of free variables \iMbox{\text{FV}(M)}} of a \stu{lambda expression}
\iMbox{M} is denoted as \iMbox{\text{FV}(M)} and \stu{defined by recursion}
on the structure of the terms, as follows:
\begin{enumerate}
    \vItem
    \iMbox{\text{FV}(x) = \{ x \}}
    \vItem
    \iMbox{\text{FV}(\lambda x. M) = \text{FV}(M) \setminus \{ x \}}
    \vItem
    \iMbox{\text{FV}(M \ N) = \text{FV}(M) \cup \text{FV}(N)}
\end{enumerate}

\kwub{Closed lambda expressions} => expression that contains \stu{no free variables},
also called \kwub{combinators}

\hSep

\kwub{Reduction} => there are three kinds:
\begin{itemize}
    \vItem \stb{$\alpha$-conversion}: changing bound variables (\stu{alpha})
    \vItem \stb{$\beta$-reduction}: applying functions to their arguments (\stu{beta})
    \vItem \stb{$\eta$-reduction}: which captures a notion of extensionality (\stu{eta})
\end{itemize}

\kwub{Equivalences} => two expressions are \stu{equivalent}, if they can be
\stu{reduced} into the \ul{same expression}
\begin{itemize}
    \vItem
    \iMbox{ M_{1} =_{\alpha} M_{2} \iff \exists M'. \ M_{1} \to_{\alpha}^{*} M' \ \text{and} \ M_{2} \to_{\alpha}^{*} M'}
    \vItem
    \iMbox{ M_{1} =_{\beta} M_{2} \iff \exists M'. \ M_{1} \to_{\beta}^{*} M' \ \text{and} \ M_{2} \to_{\beta}^{*} M'}
    \vItem
    \iMbox{ M_{1} =_{\eta} M_{2} \iff \exists M'. \ M_{1} \to_{\eta}^{*} M' \ \text{and} \ M_{2} \to_{\eta}^{*} M'}
    \vspace{1pt}

    \vItem We can \stu{combine them} to achieve \stu{more complete} notions of \stu{equivalence}:
    \begin{enumerate}
        \vItem \iMbox{ M_{1} =_{\beta \eta} M_{2} \iff \exists M'. \ M_{1}
            \to_{\beta \eta}^{*} M' \ \text{and} \ M_{2} \to_{\beta \eta}^{*} M'}
    \end{enumerate}
\end{itemize}


\hSep

\kwub{$\alpha$-conversion} => change \stu{bound variable names}:
\begin{itemize}
    \vItem Only rename variable occurrences bound by the \stu{same abstraction}

    \begin{itemize}
        \vItem
        e.g. \iMbox{ \lambda x. \lambda x. x =_{\alpha} \lambda y . \lambda x. x}
        or
        \iMbox{ \lambda x. \lambda x. x =_{\alpha} \lambda x. \lambda y. y}
        but \ul{not}
        \iMbox{ \lambda x. \lambda x. x =_{\alpha} \lambda y . \lambda x. y}
    \end{itemize}
    \vItem
    You \ul{cannot} rename a \ul{variable} if it will result in a \ul{different} \iMbox{\lambda} capturing it

    \begin{itemize}
        \vItem
        e.g. you \ul{cannot} do
        \iMbox{\lambda x. \lambda y. x =_{\alpha} \lambda y. \lambda y. y}
    \end{itemize}
\end{itemize}

\kwub{Substitution} \iMbox{E[R/x]} => replacing all \ul{free occurrences} of variable \iMbox{x} in the
expression \iMbox{E} with expression \iMbox{R}
\begin{enumerate}
    \vItem \iMbox{x} and \iMbox{y} are \textit{only} \ul{variables} \ \& \
    \iMbox{M} and \iMbox{N} are any \iMbox{\lambda}-expression
    \vItem \iMbox{\begin{aligned}
            x[\textcolor{red}{M}/y]                                                       & \equiv \begin{cases}
                                                                                                       \textcolor{red}{M} & x=y      \\
                                                                                                       x                  & x \neq y
                                                                                                   \end{cases}                                                                                                        \\
            (\textcolor{red}{M_{1}} \ \textcolor{orange}{M_{2}})[\textcolor{blue}{N} / x] & \equiv (\textcolor{red}{M_{1}}[\textcolor{blue}{N} / x]) \ (\textcolor{orange}{M_{2}} [\textcolor{blue}{N} / x])                            \\
            (\lambda x. \textcolor{red}{M})[\textcolor{orange}{N} /y]                     & \equiv \begin{cases}
                                                                                                       \lambda x. \textcolor{red}{M}                                  & x=y \vspace{2.5pt}                                                \\
                                                                                                       \lambda x. (\textcolor{red}{M}[\textcolor{orange}{N} /y])      & x\neq y \ \text{and} \ x \not\in \text{FV}(N) \vspace{2.5pt}      \\
                                                                                                       \lambda z. (\textcolor{red}{M}[z/x][\textcolor{orange}{N} /y]) & \begin{aligned}
                                                                                      & x\neq y \ \text{and} \ x \in \text{FV}(N)  \\
                                                                                      & \text{for some} \ {z \neq x,y}             \\
                                                                                      & z \not\in {\text{FV}(M) \cup \text{FV}(N)}
                                                                                 \end{aligned}
                                                                                                   \end{cases}
        \end{aligned}
    }

    \vItem To \stu{substitute} into a lambda abstraction, it is
    sometimes necessary to \stu{$\alpha$-convert} the expression:
    \begin{itemize}
        \vItem
        e.g. its \ul{wrong} for \iMbox{ (\lambda x.y)[x / y]} to
        result in \iMbox{(\lambda x.x)}, because \iMbox{x} was meant to be
        \ul{free not bound}
        \vItem
        instead the correct substitution is \iMbox{(\lambda z.x)} up to
        \stu{$\alpha$-equivalence}
    \end{itemize}
\end{enumerate}

\hSep

\kwub{$\beta$-reduction} => the idea of \stu{function application}
\begin{enumerate}
    \vItem defined in terms of \stu{substitution} where the \stu{$\beta$-reduction}
    of \iMbox{ (\lambda x. M) \ N} is \iMbox{M[N / x]}
    \begin{enumerate}
        \vItem
        e.g. \iMbox{ ((\lambda n. n \times 2) \ 7) \longrightarrow _{\beta} 7 \times 2}
    \end{enumerate}

    \vItem From there, \stu{more rules are derived}:
    \begin{enumerate}
        \vItem \begin{flushleft}
            \AxiomC{$\textcolor{red}{M} \rightarrow_{\beta} \textcolor{orange}{M'}$}
            \UnaryInfC{$\lambda x. \textcolor{red}{M} \rightarrow_{\beta} \lambda x. \textcolor{orange}{M'}$}
            \DisplayProof
            \AxiomC{$\textcolor{red}{M} \rightarrow_{\beta} \textcolor{orange}{M'}$}
            \UnaryInfC{$\textcolor{red}{M}N \rightarrow_{\beta} \textcolor{orange}{M'}N$}
            \DisplayProof
            \AxiomC{$\textcolor{red}{N} \rightarrow_{\beta} \textcolor{orange}{N'}$}
            \UnaryInfC{$M\textcolor{red}{N} \rightarrow_{\beta} M\textcolor{orange}{N'}$}
            \DisplayProof
        \end{flushleft}

        \hspace{10pt}

        \vItem \begin{flushleft}
            \AxiomC{$\textcolor{red}{M} =_{\alpha} \textcolor{red}{M'}
                \quad \quad
                \textcolor{red}{M'} \rightarrow_{\beta} \textcolor{orange}{N'}
                \quad \quad
                \textcolor{orange}{N'} =_{\alpha} \textcolor{orange}{N}$}
            \UnaryInfC{$\textcolor{red}{M} \rightarrow_{\beta} \textcolor{orange}{N}$}
            \DisplayProof
        \end{flushleft}
    \end{enumerate}

    \vItem \stb{Multiple $\beta$-reductions} \iMbox{ \to_{\beta}^{*}} =>
    \begin{enumerate}
        \vItem If we can reduce expression \iMbox{M} into expression \iMbox{M'} with
        \stu{multiple $\beta$-reductions},
        \vItem we express that fact with \iMbox{ M \to_{\beta}^{*} M'},
        \vItem with the \iMbox{*} being shorthand for the sequence of \stb{$\beta$-reductions}
    \end{enumerate}

    \vItem \stb{Reflexivity of $\alpha$-conversion} =>
    \iMbox{\frac{M={ }_\alpha M^{\prime}}{M \rightarrow_\beta^* M^{\prime}}}

    \vItem \stb{Transitivity of $\beta$-conversion} =>
    \iMbox{\frac{M \rightarrow_\beta M^{\prime \prime} \quad M^{\prime \prime}
            \rightarrow_\beta^* M^{\prime}}{M \rightarrow_\beta^* M^{\prime}}}
\end{enumerate}

\hSep

\kwub{$\eta$-reduction} => the idea of \stu{extensionality}
\begin{enumerate}
    \vItem Two functions are the \stu{same (extensionally)} \stub[orange]{iff}
    they give the \stu{same result} for \stu{all arguments}
    \vItem \begin{flushleft}
        \AxiomC{$\displaystyle x \not\in \text{FV}(\textcolor{red}{M})$}
        \UnaryInfC{$\displaystyle  (\lambda x. \textcolor{red}{M} x) \longrightarrow_{\eta} \textcolor{red}{M}$}
        \DisplayProof
    \end{flushleft}
\end{enumerate}

\hSep

\stub{Redexes and Reducts}:
\begin{enumerate}
    \vItem \kwb{reducible expression (Redex)} => \stu{subterms} that can be reduced by one of the \stu{reduction rules}
    \begin{enumerate}
        \vItem e.g. \iMbox{(\lambda x.M)} is a \stu{$beta$-redex} in expressing the
        \stu{substitution} of \iMbox{N} for \iMbox{x} in \iMbox{M}
        \vItem e.g. if \iMbox{x} is \ul{not free} in \iMbox{M}, \iMbox{\lambda x.Mx} is an \stu{$\eta$-redex}
    \end{enumerate}
    \vItem \kwb{reduct} => expression to which a \stu{redex} reduces; \\
    \stu{using the previous example}, the \stu{reducts} of these expressions are
    respectively \iMbox{M[N/x]} and \iMbox{M}

    \vItem \stub{Innermost \& Outermost $\beta$-redexes} => take the \stu{redex} \iMbox{ E = (\lambda x. M) \ N}
    \begin{enumerate}
        \vItem any \stu{redex} in \iMbox{M} or \iMbox{N} is \kwb{inside}
        of the \stu{redex} \iMbox{E}

        \vItem the \stu{redex} \iMbox{E} is \kwb{outside} of any
        \stu{redex} in \iMbox{M} or \iMbox{N}

        \vItem A \stu{redex} is \kwb{outermost} if there are no
        \textit{larger} \stu{redexes} \stu{outside} it which
        \textit{contain} it

        \vItem A \stu{redex} is \kwb{innermost} if there are no
        \stu{redexes} \stu{inside} it which it \textit{contains}

        \vItem \mkImg[100pt]{image25}
    \end{enumerate}

    \vItem \stub{Head \& Internal $\beta$-redexes} => \stu{redex} \iMbox{r} is in the \stb{Head Position} in
    expression \iMbox{E} if it has this shape:
    \begin{enumerate}
        \vItem \iMbox{ \lambda x_{0}\dots \lambda x_{n} . \textcolor{red}{(\lambda x. N) \ M} \ M_{0} \dots M_{m}},
        where \iMbox{n,m \geq 0}
        \begin{itemize}
            \vItem Here, \iMbox{ r \equiv \textcolor{red}{(\lambda x. N) \ M}}
            is the \kwb{head position} and is called the \stb{Head $\beta$-redex}
        \end{itemize}
        
        \vItem \stb{Head $\beta$-reduction} => \stu{$\beta$-reduction} applied to \stu{head position}\\
        \begin{enumerate}
            \vItem e.g. here applied to \iMbox{ r \equiv \textcolor{red}{(\lambda x. N) \ M}}
            \vItem e.g. \iMbox{\begin{aligned}
            \lambda x_{0}\dots \lambda x_{n} . &\textcolor{red}{(\lambda x. N) \ M} \ M_{0} \dots M_{m} \\
            \longrightarrow \ &\lambda x_{0}\dots \lambda x_{n} . \textcolor{red}{N[M / x]} \ M_{0} \dots M_{m}
            \end{aligned}}
        \end{enumerate}
        
        \vItem Any other \stu{$\beta$-reduction} is called an \stb{Internal $\beta$-reduction} and that 
        \stu{redex} is called an \stb{Internal $\beta$-redex}
    \end{enumerate}
\end{enumerate}

\hSep

\subsection*{Confluence (Church-Rosser property)}

The property of \textbf{\textit{confluence}} is a general \ul{rewriting
    systems} property. The idea of \textbf{\textit{confluence}} is that, if
\iMbox{a} \ul{reduces} to both \iMbox{b} and \iMbox{c} in \iMbox{\geq 0}
rewrite steps then both \ul{reducts} must in turn reduce to some
common \iMbox{d}.
\textcolor{red}{\textbf{Pasted image 20241121132641.png}}\textcolor{red}{\textbf{Pasted image 20241121133521.png}}
When applied to \ul{lambda calculus reductions}, this property is known
as the \textbf{\textit{Church--Rosser property}}. It has been stated and
proven for

\begin{itemize}

    \vItem
    \ul{Confluence} of \ul{\iMbox{\beta}-reductions}:
    \iMbox{ \forall M,M_{1},M_{2}. \ \ M \to_{\beta}^{*} M_{1} \ \text{and} \ M \to_{\beta}^{*} M_{2} \implies \exists M'. \ M_{1} \to_{\beta}^{*} M' \ \text{and} \ M_{2} \to_{\beta}^{*} M'}
    \vItem
    \ul{Confluence} of \ul{\iMbox{\eta}-reductions}:
    \iMbox{ \forall M,M_{1},M_{2}. \ \ M \to_{\eta}^{*} M_{1} \ \text{and} \ M \to_{\eta}^{*} M_{2} \implies \exists M'. \ M_{1} \to_{\eta}^{*} M' \ \text{and} \ M_{2} \to_{\eta}^{*} M'}
    \vItem
    \ul{Confluence} of \iMbox{\beta \eta}\ul{-reductions}:
    \iMbox{ \forall M,M_{1},M_{2}. \ \ M \to_{\beta \eta}^{*} M_{1} \ \text{and} \ M \to_{\beta \eta}^{*} M_{2} \implies \exists M'. \ M_{1} \to_{\beta \eta}^{*} M' \ \text{and} \ M_{2} \to_{\beta \eta}^{*} M'}
\end{itemize}

\subsection*{Normalization}

A lambda expression that \ul{cannot} be \ul{reduced further} is in
\textbf{\textit{normal form}}. You can still apply
\ul{\iMbox{\alpha}-conversion} to expressions in \textbf{\textit{normal
        form}}, in fact all \textbf{\textit{normal forms}} that can be
\textit{converted} into each other by \ul{\iMbox{\alpha}-conversion} are
defined to be \ul{equal}. Here is a summary of \textbf{\textit{normal
        forms}}:



\^{}these \textbf{\textit{normal forms}} form a ``subset'' hierarchy:
\iMbox{\beta \eta \text{-NF} \subset \beta \text{-NF} \subset \text{HNF} \subset \text{WHNF}}

Normal forms \textit{(both \ul{\iMbox{\beta \eta}-NFs} and
    \ul{\iMbox{\beta}-NFs})} are \ul{unique/equivalent} up to an
\ul{\iMbox{\alpha}-conversion},

\begin{itemize}

    \vItem
    \iMbox{\forall M,N_{1},N_{2}. \ \ M \to_{\beta \eta}^{*} N_{1} \ \text{and} \ M \to_{\beta \eta}^{*} N_{2} \ \text{and} \ N_{1},N_{2} \ \text{are in} \ \beta \eta \text{-normal form} \implies N_{1} =_{\alpha} N_{2}}
    \vItem
    \iMbox{\forall M,N_{1},N_{2}. \ \ M \to_{\beta}^{*} N_{1} \ \text{and} \ M \to_{\beta}^{*} N_{2} \ \text{and} \ N_{1},N_{2} \ \text{are in} \ \beta \text{-normal form} \implies N_{1} =_{\alpha} N_{2}}
    This is \textit{merely} a corollary to them \ul{both} having the
    \ul{Church-Rosser property},
    \vItem
    If expression \iMbox{M} has two \textbf{\textit{NFs}}
    \iMbox{N_{1},N_{2}} then \iMbox{M} \ul{must} have produced them via
    \ul{reductions} \textit{(\iMbox{\beta \eta} or \iMbox{\beta}
        respectively)}
    \vItem
    Therefore by \ul{confluence}, \iMbox{N_{1}} and \iMbox{N_{2}} must
    \ul{both} arrive at some \ul{common} \iMbox{M'} with \iMbox{\geq 0}
    \ul{reductions} *(\iMbox{\beta \eta} or \iMbox{\beta} respectively)
    \vItem
    But since \iMbox{N_{1}}, \iMbox{N_{2}} are \textbf{\textit{NFs}} they
    contain no \ul{redexes} \textit{(\iMbox{\beta \eta} or \iMbox{\beta}
        respectively)} so \iMbox{N_{1}}, \iMbox{N_{2}} arrive at \iMbox{M'}
    with exactly \iMbox{0} \ul{reductions} \textit{(\iMbox{\beta \eta} or
        \iMbox{\beta} respectively)}
    \vItem
    Meaning they were either \ul{already equal}
    \textit{(i.e. \iMbox{N_{1} \equiv M \equiv N_{2}})} or they needed to
    \ul{\iMbox{\alpha}-convert}
    \textit{(i.e. \iMbox{N_{1} =_{a} M' =_{a} N_{2}})} making them
    \ul{equivalent} up to an \ul{\iMbox{\alpha}-conversion}
    \textbf{\textit{NOTE:}} not \textit{all} expressions
    \textit{necessarily} have a \textbf{\textit{normal form}},
    e.g. \iMbox{(\lambda x. xx)(\lambda x.xx) \longrightarrow_{\beta} (\lambda x. xx)(\lambda x.xx)}
    so this expression will \textit{always} have a
    \ul{\iMbox{\beta}-redex}
\end{itemize}

\subsection*{\texorpdfstring{\iMbox{\beta \eta}-Normal
        Form}{-Normal Form}}

A term is in \textbf{\textit{\iMbox{\beta \eta}-Normal Form}} if no
\ul{\iMbox{\beta}-reductions} or \ul{\iMbox{\eta}-reductions} are
possible, or alternatively, if it contains no \ul{\iMbox{\beta}-redexes
    or \iMbox{\eta}-redexes}.

\begin{itemize}

    \vItem
    \iMbox{M \ \text{is in} \ \beta \eta \text{-normal form} \ \ \coloneq \ \ \neg \exists M'. \ M \to_{\beta} M' \ \text{or} \ M \to_{\eta} M'}
    \vItem
    \iMbox{M \ \text{has a} \ \beta \eta \text{-normal form} \ \ \coloneq \ \ \exists M'. \ M \to_{\beta \eta}^{*} M' \ \text{and} \ M' \ \text{is in} \ \beta \eta \text{-normal form}}
\end{itemize}

\subsection*{\texorpdfstring{\iMbox{\beta}-Normal Form}{-Normal Form}}

A term is in \textbf{\textit{\iMbox{\beta}-Normal Form}} if no
\ul{\iMbox{\beta}-reductions} are possible, or alternatively, if it
contains no \ul{\iMbox{\beta}-redexes}.

\begin{itemize}

    \vItem
    \iMbox{M \ \text{is in} \ \beta \text{-normal form} \ \ \coloneq \ \ \neg \exists M'. \ M \to_{\beta} M'}
    \vItem
    \iMbox{M \ \text{has a} \ \beta \text{-normal form} \ \ \coloneq \ \ \exists M'. \ M \to_{\beta}^{*} M' \ \text{and} \ M' \ \text{is in} \ \beta \text{-normal form}}
\end{itemize}

\subsection*{Head Normal Form (HNF)}

A term is in \textbf{\textit{Head Normal Form (HNF)}} if it contains no
\ul{head \iMbox{\beta}-reducts}.

\begin{itemize}

    \vItem
    Terms in \textbf{\textit{Head Normal Form (HNF)}} therefore have the
    following shape

    \begin{itemize}

        \vItem
        \iMbox{ \lambda x_{0}\dots \lambda x_{n} . \textcolor{red}{x} \ M_{0} \dots M_{m}},
        where \iMbox{x} is a \ul{variable} and \iMbox{n,m \geq 0}
    \end{itemize}
    \vItem
    \textbf{\textit{Head Normal Form (HNF)}} terms are not
    \textit{necessarily} also \ul{\iMbox{\beta \eta}-NF} or
    \ul{\iMbox{\beta}-NF}
\end{itemize}

\subsection*{Weak Head Normal Form (WHNF)}

A term is in \textbf{\textit{Weak Head Normal Form (WHNF)}} if its
either

\begin{itemize}

    \vItem
    in \ul{head normal form}, or
    \vItem
    it is a \ul{lambda abstraction} This means a WHNF term can contain
    \ul{\iMbox{\beta}-redexes} in the body of the lambda.
\end{itemize}

\subsection*{Reduction Strategies}

\textcolor{red}{\textbf{Pasted image 20241121182543.png}} \^{}quick
summary of main strategies

\subsection*{Normal Order}

\begin{itemize}

    \vItem
    \ul{\iMbox{\beta}-reduces} the \ul{leftmost + outermost redex} first

    \begin{itemize}

        \vItem
        i.e. whenever possible, arguments are substituted into the body of
        an abstraction before the arguments are reduced
    \end{itemize}
    \vItem
    \ul{Always} reduces a term to its \ul{\iMbox{\beta}-normal form}
    \textit{(if it exists)}
    \vItem
    Can perform computations in unevaluated function bodies
    \vItem
    Is not used by any programming language
\end{itemize}

\subsection*{Applicative Order}

\begin{itemize}

    \vItem
    \ul{\iMbox{\beta}-reduces} the \ul{leftmost + innermost redex} first

    \begin{itemize}

        \vItem
        therefore,  a function's arguments are always reduced before they
        are substituted into the function
    \end{itemize}
    \vItem
    Does not always reduce a term to its \ul{\iMbox{\beta}-normal form}
    \textit{(even if exists)}

    \begin{itemize}

        \vItem
        e.g. \iMbox{ ( \ \ \ \lambda x. y \ \ (\lambda z. (zz) \lambda z. (zz) \ \ \ )}
        is reduced to itself by \textbf{\textit{applicative order}},
        \vItem
        while \ul{normal order} reduces it to its \ul{\iMbox{\beta}-normal
            form} of \iMbox{y}
    \end{itemize}
\end{itemize}

\subsection*{\texorpdfstring{Full
        \ul{\iMbox{\beta}-reductions}}{Full -reductions}}

\begin{itemize}

    \vItem
    Any \ul{\iMbox{\beta}-redex} can be \ul{\iMbox{\beta}-reduced} at any
    time
    \vItem
    This means essentially the lack of any \textit{particular}
    \ul{reduction strategy}

    \begin{itemize}

        \vItem
        With regard to \ul{reducibility}, ``all bets are off''
    \end{itemize}
\end{itemize}

\subsection*{Call by Name}

\begin{itemize}

    \vItem
    Like \ul{normal order}, but \ul{no reductions are performed inside
        abstractions}

    \begin{itemize}

        \vItem
        e.g. \iMbox{\lambda x. (\lambda y.y) x} will \ul{not} be
        \ul{reduced} by this strategy
    \end{itemize}
    \vItem
    Does not always reduce a term to its \ul{\iMbox{\beta}-normal form}
    \textit{(even if exists)}
    \vItem
    Passes the function parameters unevaluated into the function body
    \vItem
    Evaluates the passed function parameter on each use
    \vItem
    Is used, with some variations, by, for example, Algol60, Haskell, R,
    and LaTeX
\end{itemize}

\subsection*{Call by Value}

\begin{itemize}

    \vItem
    Like \ul{applicative order}, but \ul{no reductions are performed
        inside abstractions}

    \begin{itemize}

        \vItem
        e.g. \iMbox{\lambda x. (\lambda y.y) x} will \ul{not} be
        \ul{reduced} by this strategy
    \end{itemize}
    \vItem
    Does not always reduce a term to its \ul{\iMbox{\beta}-normal form}
    \textit{(even if exists)}
    \vItem
    Evaluates the function parameters before passing them into the
    function body
    \vItem
    Terminates less often than \ul{call by name}, but evaluates parameters
    only once
    \vItem
    Is used, with some variations, by, for example, C, Scheme, and OCamlX
\end{itemize}

\subsection*{Call by Need}

\begin{itemize}

    \vItem
    As \ul{call by name}, but function applications that would duplicate
    terms instead name the argument

    \begin{itemize}

        \vItem
        The argument may be evaluated ``when needed''
        \vItem
        at which point the name binding is updated with the reduced value
    \end{itemize}
    \vItem
    This can save time compared to \ul{normal order} evaluation.
\end{itemize}

\subsection*{Optimal Reduction}

As \ul{normal order}, but \ul{computations} that have the \ul{same
    label} are \ul{reduced simultaneously}